\SourceLeaf{371}%
dass die Gewissheit der Geometrie nicht in der Anschauung begründet sei. Sie beruhe auf Definitionen, dem Identitätsprinzip und dem Satz vom Widerspruch. \textsc{Kant} wollte wohl nicht bestreiten, dass diese logischen Hilfsmittel ständig gebraucht werden, um den Wahrheiten der Geometrie eine Form zu geben und sie miteinander zu verknüpfen. Doch wer mit dem Wesen der Geometrie vertraut ist, kann wohl kaum verkennen, dass sie für sich allein nichts leisten und nur taube Blüten treiben, wenn die befruchtende lebendige Anschauung des Gegenstandes nicht überall wirksam ist. Der Widerspruch von Herrn \textsc{Schwab} scheint übrigens zum Teil nur auf einem Missverständnis zu beruhen. Jedenfalls scheint uns nach dem 16. Abschnitt seiner Schrift, der von Anfang bis Ende gerade das Anschauungsvermögen in Anspruch nimmt und zuletzt beweisen soll, postulata \textsc{Euclidis} in generaliora resolvi posse, non sensu et \textit{intuitione} sed \textit{intellectu} fundata, dass Herr \textsc{Schwab} mit diesen Bezeichnungen für verschiedene Zweige des Erkenntnisvermögens etwas anderes gemeint haben muss als der Königsberger Philosoph.

Obwohl der Verfasser der zweiten Schrift seinen Gegenstand ganz anders und wirklich mathematisch behandelt hat, können wir das Ergebnis doch nicht günstiger beurteilen. Wir wollen hier nicht den ganzen Gang seines Beweisversuchs darlegen und begnügen uns damit, den Punkt hervorzuheben, auf den es im Grunde ankommt. Man stelle sich zwei Geraden vor, die einander im Punkt \(N\) unter rechten Winkeln schneiden, und fälle von einem Punkt \(S\), der außerhalb dieser Geraden, aber in derselben Ebene liegt, die Senkrechten \(ST\) und \(SM\) auf die Geraden. Nun ist zu beweisen, dass \(MST\) ein rechter Winkel ist. Der Verfasser versucht einen Widerspruchsbeweis. Zuerst nimmt er an, \(MST\) sei spitz, fällt von \(T\) auf \(MS\) die Senkrechte \(Tp\) und beweist, dass \(p\) zwischen \(S\) und \(M\) liegen muss. Dann fällt er von \(p\) auf \(NT\) die Senkrechte \(pq\), wobei \(q\) zwischen \(T\) und \(N\) liegen wird. Danach fällt er von \(q\) auf \(MS\) die Senkrechte \(qp'\), wobei \(p'\) zwischen \(p\) und \(M\) liegen wird. Dann fällt er wieder von \(p'\) auf \(NT\) die Senkrechte \(p'q'\), und so weiter. Diese Schritte lassen sich ohne Ende fortsetzen. So werden von der Linie \(MS\) nacheinander die Stücke \(Sp,pp'\) und so weiter abgeschnitten, die jeweils eine angebbare Größe haben und deren Anzahl unbegrenzt ist. Der Verfasser meint nun, dass dies ein Widerspruch sei, weil \(MS\) auf diese Weise am Ende notwendig ausgeschöpft werden müsste.
